\documentclass[12pt]{article}
\usepackage[ left=0.5in,right=0.5in ]{geometry}
\usepackage{../styles/ntalgo}

\usepackage[backend=biber,style=alphabetic]{biblatex}
\addbibresource{refs.bib}

\newtheorem*{exercise}{Exercise}

\begin{document}

\section{The Framework}

We have a directed graph $G = (V, E)$ and a length function $\len : E \to \Z$.
The sizes of $V$ and $E$ are denoted $n$ and $m$, respectively. Let $U$ denote
the magnitude of the most negative weight; so $\len{u, v} \geq -U$ for all
$(u, v) \in E$. The goal is, given some $s \in V$, to compute the distance
$\dist{s, v}$ from $s$ to all $v \in V$.

\subsection{Potential Functions}

A \emph{potential function} is a weighting $\varphi : V \to \R$ on the vertices
of $G$. Potentials induce a reweighting of the graph:
\[
  \len{u, v}_\varphi \defeq \len{u, v} + \varphi(u) - \varphi(v).
\]
The reweighted graph is denoted $G_\varphi$. The length function in turn leads
to a changed distance function
\[
  \dist{s, t}_\varphi \defeq \min\Set{\sum_{e \in P} \len_\varphi(e) \mid
    \text{$P$ an $(s, t)$ walk}}.
\]
For any $(s, t)$ path $P$, the shifts of the potential function cancel out at
every intermediate vertex, so
$\len{P}_\varphi = \len{P} + \varphi(s) - \varphi(t)$, so
$\dist{s, t}_\varphi = \dist{s, t} + \varphi(s) - \varphi(t)$. Therefore,
solving SSSP on $G$ is equivalent to the same problem on $G_\varphi$.

\paragraph{Johnson's Potential}

A special class of potential functions are distance functions themselves. For
some $s \in V$ that can reach all vertices, let $\varphi(v) = \dist{s, v}$.
Then, for any $(u, v) \in E$,
\[
  \len{u, v}_\varphi = \len{u, v} + \dist{s, u} - \dist{s, v} \geq 0,
\]
since $\dist{s, v} \leq \dist{s, u} + \len{u, v}$. In fact $s$ can be replaced
by any nonempty set of vertices $S$, where
$\dist{S, v} \defeq \min_{s \in S} \dist{s, v}$. When $S = V$, $\varphi$ is the
so-called Johnson's potential.

\begin{exercise}
  Prove that $\varphi(v) \defeq - \dist{v, s}$ makes all lengths nonnegative.
\end{exercise}

Once all edges are nonnegative, we can run Dijkstra's algorithm on the
``neutralized'' graph to recover distances from any source. From now on, then,
we don't need to worry about the input source $s$.

\subsection{Scaling}

Because $G$ has integral weights, for any two paths
$P, P' : s \rightsquigarrow t$, if $\len{P} < \len{P'}$, then
$\len{P} \leq \len{P'} - 1$. Potentials don't change the difference between the
lengths of paths with common endpoints, so this remains true under any
$\varphi$. To solve SSSP, we'd like a potential $\varphi$ for which
$\len_\varphi \geq 0$ (indeed, this is what is done in the strongly polynomial
regime). Suppose, however, we can only guarantee that $\len_\varphi \geq -1/n$.
Let $\len_{\varphi, +}$ denote the length function after we round negative
lengths up to $0$. Then, if $\len{P} < \len{P'}$, since both paths have at most
$n - 1$ negative edges,
\begin{align*}
  \len{P}_{\varphi, +} <    \len{P}_\varphi + 1
                       \leq \len{P'}_\varphi
                       \leq \len{P'}_{\varphi, +},
\end{align*}
so if $P$ is a shortest path in $G_{\varphi, +}$, then it is a shortest path in
$G_\varphi$ (and thus $G$). Note that the converse is not necessarily true: if
two paths $P$ and $P'$ have equal lengths in $G$, then in $G_{\varphi, +}$, the
path with less total negative weight will win.

Of course, finding a $\varphi$ such that $\len_\varphi \geq -1/n$ seems almost
as hard as finding one that induces nonnegative lengths. We can reduce to a more
reasonable goal, though. Suppose we had a \emph{scaling algorithm} that, given a
graph with $\len \geq -U$, finds a potential $\varphi$ such that
$\len_\varphi \geq -U/2$. Successively applying this procedure $\log(nU)$ times
gives us what we want.

The rest of the work is in designing a near-linear scaling algorithm.

\section{The Algorithm}

Let's design a scaling algorithm. This is roughly based on my memory
of~\cite{BNW25}.

\subsection{SCCs Are Easy}

Recall that SSSP is trivial when $G$ is a DAG, since linear-time dynamic
programming can be used. There is a slight generalization of this: if all
negative edges travel between strongly connected components, then the problem
can be solved in $\bigO{m\log n}$ time. Let $C_1, \dots, C_k$ denote a
topological ordering of these components. For $v \in C_i$, set its potential to
$\varphi(v) = -i \cdot \text{[big number]}$. Then all negative edges $(u, v)$
have length at least $\len{u, v} + \text{[big number]} \geq 0$.

This observation also allows us to assume the input graph is strongly connected.
If it isn't, again letting $C_i$'s denote $G$'s SCCs, we can recursively find
potentials $\Set{\varphi_i}_{i \in [k]}$ that make each $C_i$ nonnegative.
Stitching these together leaves the edges between SCCs as the only possible
negative ones.

\subsection{Assumptions and Notation}

\def\half{\text{half}}

Scale all edge weights by $2/U$ so that $\len \geq -2$. Now the goal is to find
a $\varphi$ so that $\len_\varphi \geq -1$. Equivalently, letting $G_\half$
denote the graph with length function
\[
  \len{u, v}_\half
  = \begin{cases}
      \len{u, v}                  & \text{if } \len{u, v} \geq 0 \\
      \min\Set{\len{u, v} + 1, 0} & \text{otherwise}
    \end{cases},
\]
we want to find a potential that makes edges in $G_\half$ nonnegative (work out
why exactly this works). Note that for every negative $(u, v)$,
$\len{u, v} \leq 2\len{u, v}_\half$ (hence the name)\footnote{A more natural
idea is to simply halve the negative weights; that doesn't quite work, though.
Exercise: pinpoint exactly where it would break.}. Since we'll mainly be working
in this ``rounded up'' graph, from now on let $G$ and $\len$ denote $G_\half$
and $\len_\half$, with $G_0$ and $\len_0$ denoting the original graph and length
function. Let $G^+$ denote $G_\half$ with all negative lengths rounded up to
$0$.

For simplicity, we'll assume that $G_0$ has no negative cycles. We'll also
assume $G$ is strongly connected (the previous subsection shows why we can do
this).

\subsection{Bounded Positivity}

What does working in $G = G_\half$ buy us? Suppose we had a bound $\Delta$ on
the diameter of $G$, so $\dist{s, t} \leq \Delta$ for all $s, t \in V$. We want
to compute Johnson's potential $\varphi(v) \defeq \dist{V, v}$. Note that this
potential is at most $0$ (why?), so we only need to worry about paths that have
nonpositive length in $G$. Let $Q$ be any path between any $s$ and $t$ with
$\len{Q} \leq 0$. Let $P$ and $N$ denote the total positive and negative length
along $Q$, that is,
\begin{align*}
  P & = \sum_{\substack{e \in Q \\ \len{e} \geq 0}} \len{e} \\
  N & = \sum_{\substack{e \in Q \\ \len{e}<0}} \len{e}.
\end{align*}
We know that $\len{Q} = P + N \leq 0$. What about the path's length $\len_0(Q)$
in $G_0$? The negativity in $G$ is at most half that of $G_0$, so
$\len_0(Q) \leq P + 2N$. We can form a cycle by connecting $t$ to $s$ with a
path of length at most $\Delta$ (via our diameter bound; its length in $G_0$ is
no larger) to get a cycle of length at most $P + 2N + \Delta$. Since we assumed
$G_0$ has no negative cycles, $P + 2N + \Delta \geq 0$. Combining the two
inequalities gives
\[
  0 \leq -2(P + N) + P + 2N + \Delta \Longrightarrow P \leq \Delta,
\]
that is, there is a bounded amount of positivity along the path. It might not be
clear why this is useful (one may be tempted to instead bound the negativity
with $-N \leq \Delta$), but it's important enough that I'll package it in a
lemma.

\begin{lemma}
  \label{lemma:bounded-positivity}
  If $G$ has diameter at most $\Delta$ and $Q$ is some nonpositive-length path
  in $G$, then
  \[
    \sum_{e \in Q, \len{e} \geq 0} \len{e} \leq \Delta.
  \]
\end{lemma}

The main use of this lemma is in bounding the number of edges cut by a low
diameter decomposition.

\subsection{Low Diameter Decompositions}

The main tool the scaling algorithm uses is a fast algorithm for directed low
diameter decompositions (LDDs). I'll assume familiarity with this concept and
just import a lemma:

\begin{lemma}
  \label{lemma:low-diameter-decomp}
  Given a graph $G$ with nonnegative weights and a parameter $\delta$, in
  $\bigO{m \log^3 n}$ time, one can compute a randomized set of edges $S$ such
  that
  \begin{itemize}
    \item $\Pr{e \in S} \leq \bigO{\frac{\len{e}\log^2n}{\delta}}$ for each
          $e \in E$.
    \item For each SCC $C$ of $G \setminus S$ and $s, t \in C$,
          $\dist_G{s, t} \leq \delta$.
  \end{itemize}
\end{lemma}

\subsection{Lazy BFD}

This is a standard combination of Bellman-Ford and Dijkstra's algorithm
(commonly called BFD). The algorithm interleaves Bellman-Ford relaxations with
iterations of Dijkstra's algorithm:

% badness-format skip
\begin{centeralgo}
  \textsc{BFD}$(G = (V, E), \len, s)$\+
  \\  $d[v] \gets \infty$ for all $v$
  \\  $d[s] \gets 0$
  \\  $Q \gets $ priority queue on $d$ containing $s$
  \\  while $Q$ isn't empty\+
  \\    \Comment{Dijkstra's iteration}
  \\    while $Q$ isn't empty\+
  \\      Pull $u$ from $Q$
  \\      for each $(u, v) \in E$ with $\len{u, v} \geq 0$\+
  \\        if $d[u] + \len{u, v} < d[v]$\+
  \\            $d[v] \gets d[u] + \len{u, v}$
  \\            Insert $v$ into $Q$ (if not already there)\-\-\-
  \\    \Comment{Relax negative edges}
  \\    for each negative $(u, v) \in E$\+
  \\      if $d[u] + \len{u, v} < d[v]$\+
  \\        $d[v] \gets d[u] + \len{u, v}$
  \\        Insert $v$ into $Q$
\end{centeralgo}
\begin{exercise}
  Prove inductively that the value of $d[u]$ after the Dijkstra step in the
  $i$th outer iteration is the minimum distance from $s$ to $u$ using at most
  $i - 1$ negative edges.
\end{exercise}
Observe that an edge $(u, v)$ is relaxed in the $i$th iteration only if there is
some $(s, u)$-path with $i - 1$ negative edges shorter than all paths with at
most $i - 2$ negative edges. Therefore, the total number of relaxations the
algorithm performs is bounded by
\[
  \bigO{\log n\sum_{(u, v) \in E} \paren{1 + \text{min \# of negative edges on a
        shortest $(s, u)$ path}}}.
\]

\subsection{The Algorithm}

Our input is $G = G_\half$ and some $\Delta \geq 0$; we are promised that every
$s, t \in G$ has $\dist{s, t} \leq \Delta$ in the original graph (i.e.\ the
diameter is \emph{weak}). The output is a potential function
$\varphi : V \to \R$ such that every edge is nonnegative. Let $G^\star$ denote
$G$ with the addition of a supersource $s^\star$ with $0$-length edges to each
vertex. We seek to return $\dist{s^\star, \cdot} = \dist{V, \cdot} \leq 0$ as
our potentials.

We begin by cutting up the graph: with parameter $\delta = \Delta/2$, we compute
an LDD $S$ of $G^+$ in $\bigO{m \log^3 n}$ time. Note that, by the LDD
guarantee, no negative edge is in $S$ (it was rounded up to $0$). We next
compute the SCCs $C_1, \dots, C_k$ of $G \setminus S$ and recursively compute
functions $\varphi_1, \dots, \varphi_k$ that make these subgraphs nonnegative
(that is, for each $(u, v) \in E$ where $u, v \in C_i$,
$\len{u, v}_{\varphi_i} \geq 0$.) Next, we stitch these potentials together
(``SCCs are easy'') to form $\varphi$. Now $\len{u, v}_\varphi \geq 0$ for every
edge $e \notin S$.
\begin{center}
  [Insert figure of this]
\end{center}

Finally, put $\varphi(s^\star) = 0$.

Now the only remaining negative edges are in $S$. Let $u \in V$, and consider a
shortest $(s^\star, u)$ path in $G^\star$ (fixed before computing $S$). Let $Q$
denote the subpath after the first edge. In $G^\star$, this path had nonpositive
length. Since $G$ had weak diameter at most $\Delta$,
Lemma~\ref{lemma:bounded-positivity} implies that $Q$ has total positivity at
most $\Delta$. We can use this to bound the expected number of edges in
$Q \cap S$.
\begin{align*}
  \E{\sizeof{Q \cap S}} & = \sum_{e \in Q, \len{e} \geq 0}\Pr{e \in S}                          & \text{($S$ only contains positive edges)} \\
                        & = \bigO{\sum_{e \in Q, \len{e} \geq 0}\frac{\len{e}\log^2 n}{\delta}} & \text{(LDD guarantee)}                    \\
                        & = \bigO{\prac{\log^2 n}{\delta}\Delta}                                & \text{(bounded positivity)}               \\
                        & = \bigO{\log^2 n}
\end{align*}

Applying $\varphi$ doesn't change shortest paths. Therefore, along the same path
from $s^\star$ to $u$, we'd expect to see $\bigO{\log^2 n}$ negative edges (plus
one more for the initial edge). The expected running time of BFD from $s^\star$
is then
\[
  \bigO{\log n \sum_{(u, v) \in E}\paren{1 + \text{expected \# of negative edges
        in shortest $(s^\star, u)$-path}}}
  = \bigO{m\log^3 n}.
\]

The outputs of BFD are the distances $\dist{ s^\star, \cdot }_\varphi$; we
return the potentials
$\dist{s^\star, \cdot} = \dist{s^\star, \cdot}_\varphi + \varphi$.

So the nonrecursive work is $\bigO{m\log^3 n}$ for the LDD and $\bigO{m\log^3
  n}$ for the cleanup. The diameter gets halved at each level, so there are
$\bigO{\log \Delta}$ levels of recursion, giving a final running time of
$\bigO{m\log^3 n \log \Delta}$ for the scaling algorithm.

\paragraph{Hold on, what's the base case?}

When $\Delta \leq 1$, there cannot be any negative edges in $G$. If there were,
then since $\len{u, v} < 0$, $\len_0{u, v} < -1$, so we can form a
$u \to v \rightsquigarrow u$ cycle of length less than $-1 + 1 = 0$, a
contradiction. Therefore, the recursion does terminate after $\bigO{\log
  \Delta}$ levels.

\begin{nudge}
  If no negative edge shows up in any base case, where did they go? They can't
  be in the cut edges\dots
\end{nudge}

\begin{remark}
  Looking online, many of these algorithms seem to maintain this parameter
  $\kappa$, which is the maximum number of negative edges along a
  negative-weight shortest path that minimizes the number of negative edges.
  This seems very unnatural to me, since it implies that this parameter is what
  we must leverage in the recursion. It doesn't need to be, and I think the
  diameter is much more natural.
\end{remark}

\subsection{The Final Runtime}

\paragraph{Turning $\log \Delta$ into $\log n$.}

Recall that our scaling algorithm required an initial upper bound $\Delta$ on
the diameter of $G = G_\half$. Where does this come from, and why can't it be
too small? First observe that we can cap edge lengths at $2n$ in $G_0$. A
potential that neutralizes lower weights neutralizes all higher weights as well,
so the returned $\varphi$ still has $\len_\varphi \geq 0$. We also do not
introduce any negative cycles in $G_0$---a simple cycle has at most $n$ negative
edges, each with length at least $-2$, so no cycle containing an edge of length
$2n$ can be negative. Finally, since $G$ is strongly connected, we can put
$\Delta = 2n^2$ as an upper bound: every pair of vertices is connected by a path
with $n - 1$ edges of length at most $2n$. The scaling algorithm's expected
runtime is then $\bigO{m\log^4 n}$.

\paragraph{Putting it all together.}

To sum up, the algorithm successively applies scaling for $\bigO{\log(nU)}$
iterations, after which it rounds up negative weights and runs Dijkstra's
algorithm. The final (expected) runtime is
\[
  \bigO{m\log^4(n)\log(nU)}.
\]
To get a high probability guarantee, standard techniques show we just need to
pay one more log.

\section{Faster}

A reader familiar with the details of directed LDD algorithms might notice the
following. These algorithms iteratively carve out in- or out-balls from
so-called ``light'' vertices (roughly, vertices $u$ for which
$\outball{u, \delta}$ or $\inball{u, \delta}$ contains at most a constant
fraction of the graph), and recursively decompose those balls. The remaining
``core'' of ``heavy'' vertices is guaranteed to have low diameter (see the note
on LDDs for more detail). The $\log^2 n$ rather than $\log n$ loss comes from
there being logarithmic levels of recursion. As far as shortest paths are
concerned, however, reducing the size of the graph by a constant factor is
plenty of progress---we don't have to recursively decompose the smaller pieces
(in the LDD).

This means we only spend $\bigO{m\log^2 n}$ time performing the decomposition,
with a loss factor of $\bigO{\log n}$. The scaling algorithm recurses into those
smaller pieces with $\Delta$, while it recurses into the core with the smaller
$\Delta/2$. Each call reduces $\sizeof{V}\Delta$ by a constant factor, so we
still get $\bigO{\log n}$ levels of scaling recursion. The expected cleanup time
also reduces to $\bigO{m \log^2 n}$, since we cut far fewer edges. All in all,
the scaling algorithm takes $\bigO{m\log^3 n}$; the SSSP algorithm takes
$\bigO{m\log^3(n)\log(nU)}$. Another easy win is to use Thorup's integer
priority queue\footnote{The weights are not necessarily integral after scaling
by $2/U$. The scaling was just for exposition, however: working in units of
$U/2$ preserves integrality (after choosing $U$ to be a power of $2$.)} that
provides amortized $\log\log n$ time per operation to get a runtime of
\[
  \bigO{m \log^2(n)\log\log(n)\log (nU)}.
\]
This is quite a nice and compact observation about merging recursions, and only
an insane person would turn it into a 40-page paper~\cite{BCF23}.\footnote{This
is a joke.}

\begin{remark}
  In fact, the computation that is done in this sped up algorithm is almost
  exactly the same as in the old algorithm. In particular, edges are cut in the
  same way, since the recursive calls with the same $\Delta$ would begin by
  performing an LDD with $\Delta/2$, which is exactly what the old LDD did. The
  only difference is that the old algorithm cleaned up all these edges in one
  shot, while the new one does them in stages.
\end{remark}

\printbibliography
\end{document}
